\documentclass[a4paper,aps,prb,reprint,floatfix,nofootinbib,longbibliography]{revtex4-2}

\usepackage[T1]{fontenc}
\usepackage[utf8]{inputenc}
\usepackage{newtxtext,newtxmath}
\usepackage{graphicx}
\usepackage{amsmath,bm}
\usepackage{booktabs}
\usepackage{array}
\usepackage{siunitx}
\usepackage{xcolor}
\usepackage{tikz}
\usepackage{fancyhdr}
\usepackage[colorlinks=true,linkcolor=blue,citecolor=blue,urlcolor=blue]{hyperref}

\graphicspath{{../../matlab_A_deep/figures_A_deep/}}
\sisetup{scientific-notation=true}
\renewcommand{\arraystretch}{1.16}
\makeatletter
\renewcommand{\fnum@figure}{\textbf{FIG.~\thefigure}}
\renewcommand{\fnum@table}{\textbf{TABLE~\Roman{table}}}
\makeatother
\newcommand{\whuseal}{%
  \raisebox{-0.58cm}{\includegraphics[height=1.55cm]{whu_logo.png}}%
}
\setlength{\headheight}{48pt}
\setlength{\footskip}{52pt}
\fancypagestyle{firstpaperpage}{%
  \fancyhf{}%
  \fancyhead[L]{\raisebox{-0.56cm}{\rule{0.78\textwidth}{0.4pt}}}%
  \fancyhead[C]{\hspace{0.45cm}\raisebox{-0.30cm}{\normalsize\bfseries MIDTERM PAPER {\bfseries\itshape Spring Semester} (2026)}}%
  \fancyhead[R]{\whuseal}%
  \fancyfoot[C]{\small 2023300002027-\thepage}%
  \renewcommand{\headrulewidth}{0pt}%
  \renewcommand{\footrulewidth}{0pt}%
}
\fancypagestyle{mainpaperpage}{%
  \fancyhf{}%
  \fancyhead[L]{\small Yuyan-Ma}%
  \fancyhead[R]{\small MIDTERM PAPER}%
  \fancyfoot[C]{\small 2023300002027-\thepage}%
  \renewcommand{\headrulewidth}{0.4pt}%
  \renewcommand{\footrulewidth}{0pt}%
}
\definecolor{orcidgreen}{HTML}{A6CE39}
\newcommand{\orcidlink}{%
  \href{https://orcid.org/0009-0005-9074-0620}{%
    \tikz[baseline=-0.55ex,x=1ex,y=1ex]{%
      \fill[orcidgreen] (0,0) circle (0.9);
      \node[white,font=\sffamily\bfseries\tiny] at (0,0) {iD};
    }%
  }%
}

\begin{document}

\title{\large\bfseries\itshape From Boundary Quantization to Spectral-Measure Convergence:\\
Electronic Bands, Phonon Dispersions, Density of States, and Specific Heat\\
with Finite-Size Scaling in a One-Dimensional Monatomic Chain}

\author{Yuyan-Ma\,\orcidlink}
\affiliation{Wuhan Univ, Sch Phys \& Technol, Wuhan, Hubei, Peoples R China}

\begin{abstract}
This paper compares periodic and fixed boundary conditions for a one-dimensional monatomic chain.  The local second-difference operator fixes the bulk electronic and phonon dispersions, while the boundary condition fixes the finite wave-vector grid, endpoint sampling, degeneracy pattern, and spectral measure.  Periodic boundary conditions give traveling Bloch waves; fixed ends select standing waves.  Both systems still contain \(N\) normal modes, although the fixed-end phonon chain has no exact translational zero mode.  Exact diagonalization for \(N=10,20,50\), together with DOS, integrated DOS, Green-function trace, and finite-size scaling checks, shows how both spectra approach the same thermodynamic-limit measure.  The largest residual differences occur near band edges, at low temperature, and in boundary-local quantities.  A short Tamm-like example is included to separate ordinary finite-size boundary corrections from genuine localized surface states.

\vspace{0.5em}
\noindent Code\&Video: \url{https://github.com/mayuyan0121/solid_state_physics_midterm_paper}
\end{abstract}

\maketitle
\thispagestyle{firstpaperpage}
\pagestyle{mainpaperpage}

\section{Introduction}

Periodic boundary conditions are a remarkably effective idealization in solid-state physics \cite{AshcroftMermin,Kittel}.  By closing a finite chain into a ring, PBC preserves translational symmetry and makes Bloch waves, energy bands, phonon dispersions, and Brillouin-zone integrals appear in their simplest form.  Real samples are not rings, but for a macroscopic bulk observable the two ends normally contribute only an \(O(1)\) correction to an extensive quantity, or an \(O(1/N)\) correction per site.  This is the practical reason that band and phonon theory usually begins with PBC.

At nanoscale sizes, the same boundary terms can become visible \cite{Datta,Cardy}.  Fixed ends reorganize traveling waves into standing waves, change the allowed \(k\) or \(q\) grid, and alter the spacing statistics of the finite spectrum.  Low-temperature phonon heat capacity is especially sensitive because the smallest nonzero frequency scales as \(1/N\).  In one dimension, the van Hove singularity also makes band-edge sampling errors unusually easy to see.  The central question here is therefore not which boundary condition is absolutely correct, but what part of the spectrum is a bulk property and what part is a finite-size sampling effect.

The baseline model is kept deliberately simple: a one-dimensional monatomic chain, treated once as a nearest-neighbor electronic tight-binding problem and once as a harmonic phonon problem.  This choice is useful because the two spectra come from the same discrete second-difference structure.  It also prevents a common confusion.  A uniform monatomic nearest-neighbor chain has no intrinsic bulk gap, so ordinary boundary shifts should not be described as in-gap surface states.  Within this restricted model, PBC and fixed BC leave the bulk dispersion unchanged but change the finite quantization rule and therefore the finite spectral measure.

We use dimensionless parameters
\[
a=\hbar=t=M=C=1,\qquad \epsilon_0=0 .
\]
The required system sizes are \(N=10,20,50\).  Larger sizes up to \(N=1000\) are used only when a convergence trend is needed.

\section{Models and Boundary Quantization}

The electronic nearest-neighbor Hamiltonian is the standard one-orbital tight-binding model for a monatomic chain \cite{AshcroftMermin,Kittel},
\begin{equation}
H=\sum_n \epsilon_0 |n\rangle\langle n|
-t\sum_n\left(|n\rangle\langle n+1|+|n+1\rangle\langle n|\right),
\end{equation}
or, in real space,
\begin{equation}
E\psi_n=\epsilon_0\psi_n-t(\psi_{n+1}+\psi_{n-1}).
\end{equation}
The phonon equation for a monatomic harmonic chain is \cite{Kittel}
\begin{equation}
M\ddot u_n=C(u_{n+1}+u_{n-1}-2u_n).
\end{equation}
With \(u_n(t)=u_n e^{-i\omega t}\), it becomes
\begin{equation}
\omega^2 u_n=\frac{C}{M}(2u_n-u_{n+1}-u_{n-1}).
\end{equation}
Both problems have therefore reduced to eigenvalue problems for a discrete second-difference operator.  PBC places this operator on a ring.  Fixed BC places it on a finite interval and imposes nodes outside the two ends.

For electronic PBC, the Bloch form
\begin{equation}
\psi_n=N^{-1/2}e^{ikna}
\end{equation}
and \(\psi_{n+N}=\psi_n\) give
\begin{equation}
k_m=\frac{2\pi m}{Na},\qquad m=0,\ldots,N-1,
\end{equation}
and
\begin{equation}
E(k)=\epsilon_0-2t\cos(ka).
\end{equation}
Because \(E(k)=E(-k)\), the PBC spectrum has \(\pm k\) degeneracy except at special points such as \(k=0\) and, for even \(N\), \(k=\pi/a\).

For electronic fixed BC,
\begin{equation}
\psi_0=\psi_{N+1}=0,\qquad \psi_n=A\sin(kna),
\end{equation}
so
\begin{equation}
k_m=\frac{m\pi}{(N+1)a},\qquad m=1,\ldots,N,
\end{equation}
and
\begin{equation}
E_m=\epsilon_0-2t\cos\left(\frac{m\pi}{N+1}\right).
\end{equation}
Fixed BC samples only the half-Brillouin-zone standing-wave representatives.  This does not lose states: the \(N\) allowed integers \(m=1,\ldots,N\) still give \(N\) independent eigenvectors.

For phonons, PBC gives
\begin{equation}
q_m=\frac{2\pi m}{Na},\qquad
\omega(q)=2\sqrt{\frac{C}{M}}\left|\sin\frac{qa}{2}\right|.
\end{equation}
The \(q=0\) mode is the translational acoustic zero mode.  Fixed ends give
\begin{equation}
q_m=\frac{m\pi}{(N+1)a},\qquad
\omega_m=2\sqrt{\frac{C}{M}}\sin\left[\frac{m\pi}{2(N+1)}\right].
\end{equation}
Both boundary conditions have \(N\) normal modes, but a fixed-end chain has no exact translational zero mode.  In the heat-capacity calculation below, the PBC zero mode is excluded from the vibrational heat capacity.  Therefore \(C_V/(Nk_B)\) approaches \((N-1)/N\) at high temperature for PBC and \(1\) for fixed BC.  If normalized by the number of nonzero modes, both limits approach unity.

The relation between the two bases is elementary:
\begin{equation}
\sin(kna)=\frac{e^{ikna}-e^{-ikna}}{2i},\qquad
\cos(kna)=\frac{e^{ikna}+e^{-ikna}}{2}.
\end{equation}
Thus a standing wave is a particular linear combination of two opposite traveling waves.  Fixed BC selects this combination through endpoint nodes; it does not change the local bulk equation.

It is useful to separate mode counting from mode labeling.  PBC and fixed BC both diagonalize an \(N\times N\) electronic Hamiltonian and an \(N\times N\) phonon dynamical matrix.  What changes is the finite domain of the difference operator and the symmetry used to label eigenvectors.  In PBC, crystal momentum is a good quantum number and \(+k\) and \(-k\) are distinct traveling states.  In fixed BC, translation symmetry is lost and the real sine basis already combines the two directions.  This is the origin of the different degeneracy patterns in finite spectra.

\begin{figure*}[t]
\includegraphics[width=0.93\textwidth]{Fig1_dispersion_sampling.pdf}
\caption{\label{fig:dispersion}\small
Dispersion and allowed sampling.  Electronic and phonon dispersions in the first Brillouin zone \([-\pi,\pi]\).  Black solid curves are the bulk dispersions, filled circles are PBC samples, and open diamonds are fixed-BC standing-wave representatives.  Fixed BC samples only \(0<k,q<\pi/a\); the missing negative branch reflects the standing-wave basis, not a different bulk dispersion.}
\end{figure*}

\section{DOS, Spectral Measure, and Green Function}

For a finite spectrum \(\{\lambda_j\}_{j=1}^N\), I use the discrete spectral measure and normalized DOS \cite{ReedSimon}
\begin{equation}
\mu_N(\lambda)=\sum_{j=1}^{N}\delta(\lambda-\lambda_j),\qquad
\rho_N(\lambda)=\frac{1}{N}\sum_{j=1}^{N}\delta(\lambda-\lambda_j),
\end{equation}
and the integrated DOS
\begin{equation}
\mathcal N_N(\lambda)=\frac{1}{N}\#\{j:\lambda_j\le\lambda\}.
\end{equation}
For electrons the spectral variable is \(E\); for phonons it is \(\omega\).  In the thermodynamic limit,
\begin{equation}
\frac{1}{N}\sum_j f(E_j)\rightarrow
\frac{1}{2\pi}\int_{-\pi/a}^{\pi/a}f[E(k)]\,dk.
\end{equation}
For the tight-binding band \cite{AshcroftMermin,Kittel},
\begin{equation}
\rho(E)=\frac{1}{\pi\sqrt{4t^2-(E-\epsilon_0)^2}},
\quad E\in[\epsilon_0-2t,\epsilon_0+2t].
\end{equation}
The square-root divergence is the one-dimensional van Hove singularity.  It is also the reason that finite sampling errors look largest near \(E=\epsilon_0\pm 2t\).  The phonon DOS is
\begin{equation}
g(\omega)=\frac{2}{\pi\sqrt{\omega_{\max}^2-\omega^2}},
\qquad \omega_{\max}=2\sqrt{C/M},
\end{equation}
for \(0<\omega<\omega_{\max}\).

The same DOS can be obtained from the resolvent \cite{Economou},
\begin{equation}
\rho_N(E;\eta)=-\frac{1}{\pi N}
\operatorname{Im}\operatorname{Tr}(E+i\eta-H)^{-1}.
\end{equation}
For fixed BC the trace may be decomposed schematically as
\begin{equation}
\operatorname{Tr}G_{\rm fixed}(E)=
N g_{\rm bulk}(E)+g_{\rm boundary}(E)+o(1).
\end{equation}
This schematic form is the Green-function version of the boundary argument: the total trace has an extensive bulk part plus a boundary part.  After division by \(N\), the boundary contribution disappears from the normalized DOS, although an edge-site local DOS need not look bulk-like.

For a finite chain, \(\rho_N\) is a sum of delta functions and is not very informative as a raw plot.  I therefore use the Gaussian-broadened DOS
\begin{equation}
\rho_{N,\sigma}(E)=\frac{1}{N}\sum_j
\frac{1}{\sqrt{2\pi}\sigma}
\exp\!\left[-\frac{(E-E_j)^2}{2\sigma^2}\right],
\end{equation}
and the integrated DOS, which is a monotone staircase.  The latter gives a stable metric,
\begin{equation}
D_N=\sup_E|\mathcal N_N(E)-\mathcal N_\infty(E)|,
\end{equation}
analogous to a Kolmogorov distance between finite and limiting spectral measures.  This is a more stable comparison than asking two singular DOS curves to agree point by point.

For a smooth observable \(F(k)=f[E(k)]\), the fixed-BC sum admits an Euler--Maclaurin form, giving the finite-size correction structure used in scaling theory \cite{Cardy},
\begin{equation}
\frac{1}{N}\sum_{m=1}^{N}
F\!\left(\frac{m\pi}{N+1}\right)
=
\frac{1}{\pi}\int_0^\pi F(k)\,dk
+\frac{A_F}{N}
+\frac{B_F}{N^2}
+O(N^{-3}).
\end{equation}
The coefficient \(A_F\) collects the boundary-sensitive part of the finite sum: endpoint sampling, reflection phase, and finite-interval normalization.  Singular observables such as the unbroadened DOS do not follow this smooth estimate cleanly, which is why the numerical comparison below uses broadened DOS and IDOS.

The same formula gives a useful rule of thumb.  Periodic sums of smooth periodic functions often lose their endpoint correction because the interval has no physical end.  A fixed chain still has two reflecting ends, so the first boundary term can be visible at modest \(N\).  Near a band edge, where the dispersion is locally quadratic, a small mismatch in \(k\)-space can be magnified in an energy-domain plot.

Equivalently, the endpoint can be described by a reflection phase.  A standing-wave quantization condition may be written as
\begin{equation}
2kL+2\phi(k)=2\pi m,
\end{equation}
where \(\phi(k)\) is the reflection phase.  PBC instead imposes a loop phase condition.  The phase shift changes the finite counting function by \(O(1)\), and hence normalized spectral quantities by \(O(1/N)\).

\section{Numerical Results}

The finite matrices were constructed as sparse tridiagonal matrices and diagonalized exactly for the required sizes \(N=10,20,50\).  Analytic and numerical eigenvalues agree to roundoff precision, as shown in Table~\ref{tab:eigerr}.  This check is small but useful: if the numerical spectra were even slightly inconsistent with the analytic formulas, the later discussion of degeneracy and level spacing would be ambiguous.  Figure~\ref{fig:levels} shows the electronic levels, phonon frequencies, and spacing histograms.  PBC exhibits exact or near-zero spacings from the \(\pm k\) and \(\pm q\) degeneracies, whereas fixed BC generally does not.  Increasing \(N\) makes both spectra denser and moves both spectral measures toward the same continuum limit.

For electrons, the PBC Hamiltonian is a circulant tridiagonal matrix with corner hopping elements \(-t\), while the fixed-BC Hamiltonian is the open tridiagonal matrix.  For phonons, the same construction is applied to the dynamical matrix \(D=(C/M)(2I-T-T^\dagger)\), again adding corner terms only for PBC.  The eigenfrequencies are obtained from \(\omega_j=\sqrt{\max(\lambda_j,0)}\), with tiny negative roundoff values clipped to zero.  This clipping has no physical content; it only prevents a numerical zero mode from appearing as a small imaginary frequency.

\begin{table}[b]
\caption{\label{tab:eigerr}\small Maximum absolute error between analytic spectra and numerical diagonalization.}
\begin{ruledtabular}
\begin{tabular}{lccc}
kind and BC & \(N=10\) & \(N=20\) & \(N=50\)\\
\hline
electron PBC & \(1.11{\times}10^{-15}\) & \(1.46{\times}10^{-15}\) & \(2.28{\times}10^{-15}\)\\
electron fixed & \(8.88{\times}10^{-16}\) & \(9.99{\times}10^{-16}\) & \(1.55{\times}10^{-15}\)\\
phonon PBC & \(6.66{\times}10^{-16}\) & \(1.17{\times}10^{-15}\) & \(1.55{\times}10^{-15}\)\\
phonon fixed & \(4.44{\times}10^{-16}\) & \(7.77{\times}10^{-16}\) & \(9.99{\times}10^{-16}\)
\end{tabular}
\end{ruledtabular}
\end{table}

Table~\ref{tab:eigerr} separates physics from numerics.  The errors are at the \(10^{-15}\) level, so the degeneracies, the absence of a fixed-end zero mode, and the spacing differences in Fig.~\ref{fig:levels} are not artifacts of diagonalization.  They are consequences of putting the same local operator on two different finite domains.

\begin{figure*}[t]
\includegraphics[width=0.94\textwidth]{Fig2_levels_spacing.pdf}
\caption{\label{fig:levels}\small
Energy/frequency levels and level-spacing statistics.  (a) Electronic levels and (b) phonon frequencies for \(N=10,20,50\), grouped as \(N\) PBC and \(N\) fixed.  (c,d) Probability-normalized spacing histograms for \(\Delta E\) and \(\Delta\omega\).  The zero-spacing peak in PBC reflects \(\pm k\) or \(\pm q\) degeneracy, except at special points; fixed BC generally removes this explicit twofold degeneracy by forming standing waves.}
\end{figure*}

The top panels of Fig.~\ref{fig:levels} are the energy-level-versus-\(N\) comparison.  Each vertical strip is one finite system.  For \(N=10\) the discreteness is obvious; by \(N=50\) the same points already begin to outline a band.  The bottom panels compress this information into spacing statistics.  PBC has many zero or near-zero spacings because \(E(k)=E(-k)\) and \(\omega(q)=\omega(-q)\).  Fixed BC does not display this explicit pair structure because the eigenvectors have already combined the two traveling directions into one standing wave.

Figure~\ref{fig:dos} compares broadened DOS, integrated DOS, and the IDOS convergence metric
\begin{equation}
D_N=\sup_E|\mathcal N_N(E)-\mathcal N_\infty(E)|.
\end{equation}
The broadened DOS uses \(\sigma/t=0.08\).  The analytic van Hove peaks are clipped in the figure only for readability.  The IDOS distance decreases approximately as a boundary-order correction.  This is the main reason for plotting IDOS in addition to DOS.  A broadened DOS can look different when the smoothing width is changed, but the staircase IDOS asks a simpler question: how many states have been accumulated below a given energy?

Table~\ref{tab:doscheck} reports the normalization error of the broadened DOS.  The small residual error comes mainly from the finite energy grid and from Gaussian tails outside the plotted energy window; the total number of states is otherwise conserved.  The IDOS distance scales roughly as \(1/N\), as one would expect from the finite spacing of the allowed wave-vector grid.  Near the band edges, however, the ordinary DOS plot looks less well behaved because \(dE/dk\) vanishes and many nearby \(k\)-values are squeezed into a narrow energy interval.

\begin{table}[b]
\caption{\label{tab:doscheck}\small Electronic DOS/IDOS validation for \(\sigma/t=0.08\).}
\begin{ruledtabular}
\begin{tabular}{lccc}
BC & \(N\) & DOS norm. error & IDOS KS distance\\
\hline
PBC & 10 & \(-5.74{\times}10^{-8}\) & 0.099994\\
fixed & 10 & \(-1.83{\times}10^{-10}\) & 0.090334\\
PBC & 20 & \(-2.88{\times}10^{-8}\) & 0.050000\\
fixed & 20 & \(-6.55{\times}10^{-9}\) & 0.047256\\
PBC & 50 & \(-1.99{\times}10^{-8}\) & 0.019994\\
fixed & 50 & \(-1.45{\times}10^{-8}\) & 0.018777
\end{tabular}
\end{ruledtabular}
\end{table}

\begin{figure*}[t]
\includegraphics[width=0.95\textwidth]{Fig3_dos_idos.pdf}
\caption{\label{fig:dos}\small
DOS, IDOS, and spectral-measure convergence.  (a) Gaussian-broadened electronic DOS for \(N=50\) with \(\sigma/t=0.08\), compared with the thermodynamic-limit DOS.  The analytic van Hove singularities are clipped for readability.  (b) Integrated DOS staircase and continuous limit.  (c) \(D_N=\sup_E|\mathcal N_N(E)-\mathcal N_\infty(E)|\) versus \(N\) on log-log axes, demonstrating spectral-measure convergence.}
\end{figure*}

Figure~\ref{fig:cv} shows the phonon heat capacity.  This observable is a useful sanity check because it turns the spectral difference into a measurable thermodynamic effect.  The low-temperature part is cut off by the smallest nonzero phonon frequency: when \(k_BT\ll \hbar\omega_{\min}\), no vibrational mode is thermally populated and \(C_V\) is exponentially suppressed.  Replotting the data against \(T/\omega_{\min}\) collapses the low-temperature onset reasonably well, confirming that the main finite-size scale is \(\omega_{\min}\sim 1/N\).

The heat-capacity comparison also makes the zero-mode convention visible.  In a finite PBC chain the \(q=0\) mode corresponds to rigid translation and carries no oscillator heat capacity.  Consequently, the high-temperature limit of \(C_V/(Nk_B)\) is slightly below unity for PBC, by exactly one mode out of \(N\).  This distinction disappears in the thermodynamic limit, but for \(N=10\) it is large enough to see.  Fixed BC has no such translational degree of freedom because the endpoints are pinned.

\begin{figure*}[t]
\includegraphics[width=0.91\textwidth]{Fig4_cv_collapse.pdf}
\caption{\label{fig:cv}\small
Specific heat and finite-size collapse.  (a) \(C_V/(Nk_B)\) for finite phonon chains.  The PBC translational zero mode is excluded from vibrational heat capacity, so its high-temperature per-site limit is \((N-1)/N\); fixed BC tends to 1.  If normalized by the number of nonzero modes, both high-temperature limits tend to 1.  (b) Low-temperature collapse versus \(T/\omega_{\min}\), showing that the finite-size cutoff is controlled by the smallest nonzero phonon frequency.}
\end{figure*}

Figure~\ref{fig:green} gives the Green-function reconstruction of the DOS for fixed BC \cite{Economou}.  The broadening parameter \(\eta/t\) is logarithmically spaced from \(10^{-3}\) to \(0.2\), and the color scale is clipped at 2 to make both sharp and broad features visible.  At the smallest \(\eta\), the map resolves individual finite-chain eigenvalues as narrow Lorentzian ridges.  As \(\eta\) increases, the ridges merge into the same smooth envelope obtained by explicit DOS broadening.  This second route to the DOS is helpful because it does not rely on histogram bins.

\begin{figure}[t]
\includegraphics[width=\linewidth]{Fig5_green_map.pdf}
\caption{\label{fig:green}\small
Green-function reconstruction of broadened DOS.  The fixed-BC electronic DOS for \(N=50\) is obtained from
\(\rho_N(E;\eta)=-(\pi N)^{-1}\operatorname{Im}\operatorname{Tr}(E+i\eta-H)^{-1}\).
Increasing \(\eta\) converts discrete spectral lines into a smooth broadened DOS.  The color scale is clipped at \(\rho(E;\eta)t=2\) for visibility.}
\end{figure}

\section{Thermodynamic Limit and Surface-State Boundary}

The allowed-grid spacings are
\begin{equation}
\Delta k_{\rm PBC}=\frac{2\pi}{Na},\qquad
\Delta k_{\rm fixed}=\frac{\pi}{(N+1)a}.
\end{equation}
Both are \(O(1/N)\), so both give Riemann sums for the same bulk Brillouin-zone integral.  The bulk dispersion is fixed by the local difference equation; the boundary only changes \(O(1)\) endpoint conditions or reflection phases.  Consequently, normalized DOS, energy per site, and bulk thermodynamic observables have the same thermodynamic limit.  The most visible finite-size deviations occur near band edges, at low temperature, and in boundary-local probes.

A representative smooth observable is the half-filled spinless electronic energy per site,
\begin{equation}
e_N=\frac{1}{N}\sum_{j=1}^{\lfloor N/2\rfloor}E_j .
\end{equation}
For the present band its thermodynamic limit is
\begin{equation}
e_\infty=\frac{1}{2\pi}\int_{-\pi/2}^{\pi/2}(-2t\cos k)\,dk
=-\frac{2t}{\pi}.
\end{equation}
Numerically, \(|e_N-e_\infty|\) follows a power-law finite-size trend, with fixed BC showing the clearest boundary-order contribution.  This behavior is consistent with the Euler--Maclaurin picture: smooth integrated observables converge more regularly than singular DOS-level observables.

A true surface state must be distinguished from this ordinary boundary correction \cite{Davison,Tamm,Shockley}.  The uniform monatomic nearest-neighbor chain has a single band and no bulk gap, so a fixed end alone should not be called an in-gap surface state.  For a semi-infinite chain with a surface onsite perturbation \(\Delta\), set
\begin{equation}
\psi_n=A\lambda^{n-1},\qquad |\lambda|<1.
\end{equation}
The bulk recursion and surface equation are
\begin{equation}
E-\epsilon_0=-t(\lambda+\lambda^{-1}),\qquad
E-\epsilon_0-\Delta=-t\lambda.
\end{equation}
Thus
\begin{equation}
\lambda=-\frac{t}{\Delta},\qquad |\Delta|>|t|,
\end{equation}
and
\begin{equation}
E_b=\epsilon_0+\Delta+\frac{t^2}{\Delta}.
\end{equation}
For \(\epsilon_0=0\), \(t=1\), and \(\Delta=2\), one obtains \(\lambda=-1/2\) and \(E_b=2.5\).  Since the bulk band is \([-2,2]\), this is a bound state outside the upper band, not an in-gap state.  Genuine in-gap states require a surface potential plus a gap, hopping reconstruction, dimerization, or another structural mechanism.

The finite-chain check in Table~\ref{tab:surface} uses an open chain of \(N=100\) sites with the same surface onsite perturbation at the left endpoint.  The highest eigenvalue is already indistinguishable from the semi-infinite prediction at the displayed precision, and the boundary weight is concentrated on the first site.  This example is included only to clarify the conceptual boundary between an \(O(1/N)\) finite-size spectral correction and a genuine exponentially localized eigenstate.

\begin{widetext}
\begingroup
\refstepcounter{table}\label{tab:surface}
\begin{center}
\small
\setlength{\tabcolsep}{8pt}
\textbf{TABLE~\Roman{table}.} Tamm-like surface-state example.  The finite-chain check uses \(N=100\).
\vspace{0.5em}

\begin{ruledtabular}
\begin{tabular}{lcccccc}
model & \(\Delta/t\) & analytic \(E_b/t\) & numerical \(E_b/t\) & \(|\lambda|\) & \(\xi/a\) & boundary weight\\
\hline
surface onsite perturbation & 2 & 2.5 & 2.5 & 0.5 & 1.4427 & 0.75
\end{tabular}
\end{ruledtabular}
\end{center}
\endgroup
\end{widetext}

\section{Discussion}

\textit{Why do most textbooks use PBC?}  PBC preserves translational symmetry, makes the eigenstates Bloch waves, and converts finite sums into Brillouin-zone integrals.  It removes endpoint scattering and is therefore the cleanest route to bulk bands, phonons, DOS, and thermodynamic quantities.  In addition, PBC avoids choosing an arbitrary surface termination when the intended question concerns the bulk crystal rather than a finite sample.

\textit{Can fixed BC describe bulk properties?}  Yes, if the system is large, the observable is bulk averaged, and the probe is not explicitly surface-sensitive.  Fixed BC introduces standing waves, endpoint nodes, and the absence of a strict translational phonon zero mode, but these differences vanish as boundary-order corrections in normalized bulk observables.  Fixed BC becomes more appropriate when the experiment resolves confinement, end scattering, or local edge response.

\textit{How do nanocrystals, quantum dots, and nanowires differ from bulk bands?}  Their spectra are discrete, and the level spacing may be comparable with \(k_BT\), optical linewidths, or disorder broadening \cite{Datta}.  PBC is useful as a bulk reference; open or fixed BC is more appropriate for confinement, endpoints, and surface-sensitive physics.  In a real nanostructure one would also add passivation, reconstruction, disorder, and electrode self-energies when those effects are relevant.

\textit{How can surface states affect absorption, catalysis, and conductivity?}  They can create subgap optical channels and recombination paths, alter surface local DOS and adsorption strength in catalysis, and either provide conducting channels or act as traps and scattering centers depending on their localization, filling, and coupling to leads \cite{Davison,Datta}.  The sign of the effect is therefore material- and termination-dependent: a surface state can enhance transport if it is extended along the surface, or suppress transport if it mainly traps carriers.

\section{Conclusion}

PBC and fixed BC in a one-dimensional monatomic chain mainly change finite quantization, endpoint sampling, degeneracy, and spectral measure; they do not change the bulk dispersion.  PBC produces traveling Bloch waves and \(\pm k\) degeneracies, whereas fixed BC produces standing waves with the same number of degrees of freedom but usually without explicit twofold degeneracy.  The phonon PBC zero mode is a translational mode and is excluded from vibrational heat capacity.  Numerical spectra, DOS, IDOS, Green-function reconstruction, and heat capacity all support convergence to the same thermodynamic-limit bulk measure.  Boundary corrections are typically \(O(1/N)\) in normalized quantities, but they are amplified near band edges, at low temperature, and in local boundary responses.  The Tamm-like example clarifies that ordinary finite-size boundary corrections are not the same as true localized surface states.

\clearpage
\onecolumngrid
\appendix

\section{Supplementary boundary-error surface}

The supplementary surface below gives a visual diagnostic for the finite-size boundary correction discussed in Secs.~IV and V.  The plotted levels are first sorted within each boundary condition and then compared at the same normalized level index.  This sorting removes trivial label differences between PBC traveling-wave states and fixed-BC standing-wave states, leaving a direct view of how much the finite spectrum is shifted by the boundary.

\begingroup
\refstepcounter{figure}\label{fig:suppsurface}
\begin{center}
\includegraphics[width=0.76\textwidth,height=0.42\textheight,keepaspectratio]{Supp_boundary_error_surface.pdf}

\vspace{0.4em}
\parbox{0.84\textwidth}{\small\textbf{FIG.~\thefigure.} Supplementary boundary-error surface.  The plotted quantity is the absolute difference between sorted electronic energy levels under fixed BC and PBC, interpolated against normalized level index for the dimensionless tight-binding chain.  The surface is a diagnostic of finite-size boundary corrections rather than evidence for a new surface band.}
\end{center}
\endgroup

The surface should be read together with Fig.~\ref{fig:dos}.  In the middle of the band, where the group velocity is finite, adjacent \(k\)-points map smoothly into adjacent energies and the boundary-induced difference decreases regularly with \(N\).  Near the band edges, the dispersion is flat and the one-dimensional DOS is singular, so a comparable shift in the allowed \(k\)-grid produces a larger visible energy-domain error.  This is why boundary corrections may be small in normalized integrated quantities but still noticeable in band-edge spectroscopy or low-temperature observables.

All plotted quantities are dimensionless.  The required sizes \(N=10,20,50\) are used in the main spectra and spacing figures; larger sizes enter only the convergence checks.

\twocolumngrid

\begin{thebibliography}{10}

\bibitem{AshcroftMermin}
N. W. Ashcroft and N. D. Mermin, \textit{Solid State Physics} (Holt, Rinehart and Winston, New York, 1976).

\bibitem{Kittel}
C. Kittel, \textit{Introduction to Solid State Physics}, 8th ed. (Wiley, Hoboken, 2005).

\bibitem{Datta}
S. Datta, \textit{Quantum Transport: Atom to Transistor} (Cambridge University Press, Cambridge, 2005).

\bibitem{Cardy}
J. L. Cardy, ed., \textit{Finite-Size Scaling} (North-Holland, Amsterdam, 1988).

\bibitem{ReedSimon}
M. Reed and B. Simon, \textit{Methods of Modern Mathematical Physics} (Academic Press, New York, 1972--1979).

\bibitem{Economou}
E. N. Economou, \textit{Green's Functions in Quantum Physics}, 3rd ed. (Springer, Berlin, 2006).

\bibitem{Davison}
S. G. Davison and M. St\k{e}\'slicka, \textit{Basic Theory of Surface States} (Oxford University Press, Oxford, 1992).

\bibitem{Tamm}
I. E. Tamm, Phys. Z. Sowjetunion \textbf{1}, 733 (1932), commonly cited original Tamm surface-state work.

\bibitem{Shockley}
W. Shockley, Phys. Rev. \textbf{56}, 317 (1939).

\end{thebibliography}

\end{document}
















